% latex/supplement/figures.tex
%
% Figure captions for the seven data figures the extended suite
% reproduces. Each caption records what the figure shows, what the
% underlying data is, and the specific finding it supports. The
% captions are the reference; the figure itself is generated by
% extended_experimentation/reporting/a3e_figures.py.

\section{Figure captions}

The seven figures reproduced by
\texttt{extended\_experimentation/reporting/a3e\_figures.py} are
described below. A reviewer regenerating them from the pipeline
should see the same qualitative shapes the thesis reports; small
numerical differences in the plotted points are expected and are
covered by the tolerance bands in the evaluation suite.

\subsection*{Figure F1 --- Detection accuracy vs.\ attack severity}

\emph{Caption.} Detection accuracy across the ten-scenario attack
catalogue, plotted against severity. Each point is one scenario at one
severity variant; the horizontal dashed line marks the
$99.94\%$ thesis reference. Accuracy remains above $99.4\%$ for all
scenarios up to severity $0.65$; the adaptive-blend variant of the
tunnel scenario (S2) at severity $0.65$ is the first to fall below the
reference, at $99.31\%$.

\emph{Supports.} Section~6.2.2, Table~30.

\subsection*{Figure F2 --- Byzantine robustness vs.\ Byzantine fraction}

\emph{Caption.} Mean and maximum attack-to-reference ratios across
the Byzantine sweep. The mean ratio (blue) rises from $0.42$ at
fraction $0.05$ to $1.87$ at fraction $0.40$; the maximum ratio
(orange) rises from $0.58$ to $2.94$. The gap between the two curves
is the heterogeneity across the eight attack strategies; the mean is
what the aggregate-accuracy panel of Table~31 reports, and the
maximum is what the worst-case robustness discussion in
Section~6.3.4 cites.

\emph{Supports.} Section~6.3, Table~31.

\subsection*{Figure F3 --- DP epsilon vs.\ utility}

\emph{Caption.} Detection accuracy (green) and false-positive rate
(purple) as a function of the differential-privacy budget $\eps$. The
vertical dashed line marks the thesis's $\eps = 0.68$. At
$\eps = 0.68$ the accuracy loss relative to the non-private baseline
is $0.4\%$ and the FPR loss is $0.05\%$; both are within the
tolerances recorded in the evaluation suite. The sharp utility drop
below $\eps = 0.4$ is the reason the thesis chooses $\eps = 0.68$
rather than a smaller budget.

\emph{Supports.} Section~6.6.2, Table~32.

\subsection*{Figure F4 --- Latency CDF by protocol}

\emph{Caption.} Latency percentiles across DoH, DoT and DoQ. The
curves plot the $P_{50}$, $P_{95}$ and $P_{99}$ points of each
protocol's latency distribution. DoQ is the lowest-latency transport
at $P_{50}$ ($0.81\,\mathrm{ms}$) because its UDP-based transport
avoids the TCP handshake; DoH is the highest ($1.12\,\mathrm{ms}$) for
the same reason. The horizontal dashed line marks the thesis's
$1.42\,\mathrm{ms}$ P99 reference, which the pooled distribution
meets.

\emph{Supports.} Section~6.7.1, Table~35.

\subsection*{Figure F5 --- Throughput vs.\ concurrency}

\emph{Caption.} Peak QPS as a function of concurrency. Throughput
rises linearly up to concurrency $= 32$ and flattens thereafter as the
pipeline's inference path saturates. The horizontal dashed line
marks the $22.5\,\mathrm{M}\,\mathrm{QPS}$ thesis reference, which
corresponds to the concurrency $= 32$ point.

\emph{Supports.} Section~6.7.2, Table~35.

\subsection*{Figure F6 --- QHE NTT vs.\ convolution speedup}

\emph{Caption.} Median transform time for the negacyclic NTT
(blue) versus naive polynomial multiplication (orange). The
annotated speedup factor is the ratio of the two medians; on the
reference machine the NTT is $18.4\times$ faster. The naive
implementation is the C reference from
\texttt{c/reference/src/ntt.c}; the NTT is the Rust implementation
from \texttt{rust/a3e\_core/src/ntt.rs}.

\emph{Supports.} Section~6.5.1, Table~33.

\subsection*{Figure F7 --- Multi-protocol accuracy}

\emph{Caption.} Detection accuracy for DoH, DoT and DoQ under the
same detector. All three protocols exceed the $99.94\%$ thesis
reference; the difference between them is below the tolerance band,
which is the evidence for the claim that the detector generalises
across transports.

\emph{Supports.} Section~6.7.3, Table~37.
